\thispagestyle{empty}

\begin{tikzpicture}[
    remember picture,
    overlay,
    shift={(current page.south west)},
  ]
  % BME logo at the bottom
  \node[
    anchor=south west,
    fill=primaryColor,
    minimum width=\paperwidth,
    minimum height=2.2cm,
  ] at (0,0) {
    \includegraphics[height=1.4cm]{../../shared/bme-logo-white.pdf}
  };

  % Title
  \node[primaryColor, inner sep = 5mm] (subject) at (\paperwidth*0.5, \paperheight*0.745) {
    \sffamily\bfseries\Huge
    Matematika G1
  };
  \node[below, primaryColor] (code) at (subject.south) {
    \sffamily\LARGE
    (BMETE94BG01)
  };
  \node[below, primaryColor, inner sep=5mm] (title) at (code.south) {
    \sffamily\bfseries\LARGE
    Előadásjegyzet
  };

  % Title boxed
  \begin{pgfonlayer}{bg}
    \draw[
      primaryColor,
      line width=1mm,
      fill=primaryColor!5,
    ] (0,0) rectangle (\paperwidth, \paperheight);
    \draw[
      primaryColor,
      ultra thick,
      rounded corners,
      fill=white
    ] (subject.north west) rectangle (title.south east -| subject.north east);
  \end{pgfonlayer}

  % Authors
  \node[primaryColor, minimum width=5cm, align=center,anchor=north] (authors) at (\paperwidth*0.285, \paperheight*0.355) {
    \sffamily\itshape\Large
    Készítették: \\[3mm]
    \sffamily\bfseries\Large
    Kun László Ákos \\[2mm]
    \sffamily\bfseries\Large
    Sándor Tibor
  };
  \node[primaryColor, minimum width=5cm, align=center,anchor=north] (authors) at (\paperwidth*0.715, \paperheight*0.355) {
    \sffamily\itshape\Large
    Lektorálta: \\[3mm]
    \sffamily\bfseries\Large
    Dr. Szilágyi Brigitta
  };

  % QR codes
  \node[black, anchor=south west] (QR1) at (0.5,2.5) {
    \qrcode[height=1in]{https://github.com/szilagyi-math/notes}
  };
  \node[anchor=south west] (QT1) at (QR1.south east) {
    \sffamily%
    (1)
  };
  \node[black, anchor=south west] (QR2) at (QT1.south east) {
    \qrcode[height=1in]{https://szilagyi-math.hu/subjects/G1}
  };
  \node[anchor=south west] (QT2) at (QR2.south east) {
    \sffamily%
    (2)
  };

  \node[black, anchor=south west, minimum width=10cm, align=left] (QR3) at (QT2.south east) {
    \sffamily%
    (1) A jegyzet forráskódja\\[2mm]
    \phantom{(2)}\; (\url{https://github.com/szilagyi-math/notes})\\[2mm]
    \sffamily%
    (2) Tantárgy honlapja\\[2mm]
    \phantom{(2)}\; (\url{https://szilagyi-math.hu/subjects/G1})
  };

  % ~~~~~~~~~~~~~~~~~~~~~~~~~~~~~~~~~~~~~~~~~~~~~~~~~~~~~~~~~~~~~~~~~~~~~~~~~~~~
  % Random set operations
  % ~~~~~~~~~~~~~~~~~~~~~~~~~~~~~~~~~~~~~~~~~~~~~~~~~~~~~~~~~~~~~~~~~~~~~~~~~~~~

  % Red: A * B
  \begin{scope}[shift={(\paperwidth*.1,\paperheight*.925)}, ultra thick]
    \draw[primaryColor] (0,0) circle (.75cm);
    \draw[primaryColor] (0.75,0.5) circle (.75cm);
    \clip[] (0,0) circle (.75cm);
    \fill[primaryColor!50] (0.75,0.5) circle (0.75cm);

    \draw[primaryColor] (0,0) circle (.75cm);
    \draw[primaryColor] (0.75,0.5) circle (.75cm);
  \end{scope}

  % Blue: B - A
  \begin{scope}[shift={(\paperwidth*.1,\paperheight*.825)}, ultra thick]
    \draw[secondaryColor, fill=secondaryColor!50] (0.75,0.5) circle (.75cm);
    \draw[secondaryColor] (0,0) circle (.75cm);
  \end{scope}

  % Yellow: A - B
  \begin{scope}[shift={(\paperwidth*.25,\paperheight*.9125)}, ultra thick]
    \draw[ternaryColor, fill=ternaryColor!50] (0,0) circle (.75cm);
    \draw[ternaryColor] (0.75,0.5) circle (.75cm);
  \end{scope}

  % Green: (A - B) + (B - A)
  \begin{scope}[shift={(\paperwidth*.25,\paperheight*.8125)}, ultra thick]
    \draw[green!40!black, fill=green!40!black!50] (0,0) circle (.75cm);
    \draw[green!40!black, fill=green!40!black!50] (0.75,0.5) circle (.75cm);
    \clip (0,0) circle (.75cm);
    \fill[white] (0.75,0.5) circle (0.75cm);
  \end{scope}

  \begin{scope}[shift={(\paperwidth*.25,\paperheight*.8125)}, ultra thick]
    \draw[green!40!black] (0,0) circle (.75cm);
    \draw[green!40!black] (0.75,0.5) circle (.75cm);
  \end{scope}

  % ~~~~~~~~~~~~~~~~~~~~~~~~~~~~~~~~~~~~~~~~~~~~~~~~~~~~~~~~~~~~~~~~~~~~~~~~~~~~
  % Complex numbers
  % ~~~~~~~~~~~~~~~~~~~~~~~~~~~~~~~~~~~~~~~~~~~~~~~~~~~~~~~~~~~~~~~~~~~~~~~~~~~~

  % cartesian
  \begin{scope}[ultra thick, shift={(\paperwidth*.825,\paperheight*.865)}, rotate=60]
    \draw[->] (-.5, 0) -- (3, 0) node[above left] {$\Re$};
    \draw[->] (0, -.5) -- (0, 3) node[above left] {$\Im$};

    \draw[draw=gray, dashed]
    (1.5,0) node[below right] {$a$}
    -- (1.5, 2.5)
    -- (0, 2.5) node[below left] {$b$};

    \draw[fill=primaryColor] (1.5, 2.5) circle (0.1)
    node[above=2mm] {$z = a + b\iu$};
  \end{scope}

  % polar
  \begin{scope}[shift={(\paperwidth*.5,\paperheight*.885)}, scale=3/4]
    \foreach \ang in {0,...,31} {
        \draw [lightgray] (0,0) -- (\ang * 180 / 16:3.1);
      }

    % Concentric circles
    \foreach \s in {0, 1, 2} {
        \draw [lightgray, thick] (0,0) circle (\s + 0.5);
        \draw [primaryColor, thick] (0,0) circle (\s);
      }

    % Add the labels at multiples of pi/4
    \foreach \ang/\lab/\dir in {
    % 0/0/right,
    1/{\pi/4}/{above right},
    2/{\pi/2}/above,
    3/{3\pi/4}/{above left},
    4/{\pi}/left,
    5/{5\pi/4}/{below left},
    7/{7\pi/4}/{below right},
    6/{3\pi/2}/below} {
    \draw (0,0) -- (\ang * 180 / 4:3.25);
    }

    % Outer circle
    \draw [thick] (0,0) circle (3);

    % 0 angle line
    \draw [-to, ultra thick] (0,0) -- (3.75,0) node [above left] {$r$};

    % Example
    \draw[secondaryColor, ultra thick] (0,0) -- (60:2.75);
    \draw[secondaryColor, thick] (0,0) -- ++(150:.5);
    \draw[secondaryColor, thick] (60:2.75) -- ++(150:.5);

    \draw[draw=primaryColor, ultra thick, to-to] (150:.35) -- ++(60:2.75)
    node[midway, above, fill=primaryColor!5, rotate=60, inner sep=2pt, outer sep=2pt]
    {$r$};

    \draw[ultra thick, primaryColor, -to] (1.4,0) arc
      [start angle=0, end angle=60, radius=1.4cm];

    \node[fill=primaryColor!5, inner sep=2pt, outer sep=2pt] at (30:.85) {$\varphi$};

    \draw[fill=primaryColor] (60:2.75) circle (0.1)
    node[anchor=290, yshift=2mm, fill=primaryColor!5, inner sep=2pt, outer sep=5pt]
      {$z = r(\cos \varphi + \iu \sin \varphi)$};

  \end{scope}

  % ~~~~~~~~~~~~~~~~~~~~~~~~~~~~~~~~~~~~~~~~~~~~~~~~~~~~~~~~~~~~~~~~~~~~~~~~~~~~
  % Functions (all)
  % ~~~~~~~~~~~~~~~~~~~~~~~~~~~~~~~~~~~~~~~~~~~~~~~~~~~~~~~~~~~~~~~~~~~~~~~~~~~~

  % Sine function for integral
  \begin{scope}[shift={(\paperwidth*.1,\paperheight*.525)}]
    \draw[opacity=0, name path=r] (-0.5,0) -- (16.5,0);

    \draw[-to, ultra thick, draw=secondaryColor] (-0.75,0) -- (\paperwidth*.80,0) node[above left] {$x$};
    \draw[-to, ultra thick, draw=secondaryColor] (0,-2) -- (0,2.25) node[below right] {$y$};

    \draw[ultra thick, primaryColor, name path=s]
    plot[domain=-0.5:16.5, smooth, samples=150]
    (\x, {1.5*sin(\x r)});

    % Fill between the x-axis and the function
    \tikzfillbetween[
      of=r and s,
      split,
      every odd segment/.style={primaryColor!25},
      every even segment/.style={secondaryColor!25},
      on layer=bg,
    ]{}

    % + at pi/2 + 2k*pi
    \foreach \x in {1.571, 7.853, 14.137} {
        \node[
          circle,
          fill=primaryColor!30,
          draw=primaryColor,
          minimum size=.65cm,
          ultra thick,
        ] (A1) at (\x,0.65) {};
        \draw[primaryColor, very thick, line cap = round] ($(A1) - (1.5mm,0)$) -- ++(3mm,0);
        \draw[primaryColor, very thick, line cap = round] ($(A1) - (0,1.5mm)$) -- ++(0,3mm);
      }

    % - at 3pi/2 + 2k*pi
    \foreach \x in {4.712, 10.995} {
        \node[
          circle,
          fill=secondaryColor!30,
          draw=secondaryColor,
          minimum size=.65cm,
          ultra thick,
        ] (A2) at (\x,-0.65) {};
        \draw[secondaryColor, very thick, line cap = round] ($(A2) - (1.5mm,0)$) -- ++(3mm,0);
      }
  \end{scope}

  % -tanh(x) for fn analysis
  \begin{scope}[shift={(.5\paperwidth,.3275\paperheight)}]
    \draw[-to, draw=secondaryColor, ultra thick] (0,-3) -- (0,3.25) node[below right] {$y$};
    \draw[-to, draw=secondaryColor, ultra thick] (-8.5,0) -- (8.5,0) node[above left] {$x$};

    \draw[primaryColor, ultra thick, primaryColor!5, opacity=.75] (-7.25,0) -- ++(5.5,0);
    \draw[primaryColor, ultra thick, primaryColor!5, opacity=.75] (7.25,0) -- ++(-5.5,0);

    \draw[smooth, samples=150, ultra thick, primaryColor]
    plot[domain=-6.8:6.8] ({1.25*\x}, {-2.5*tanh(\x)});

    \draw[dashed, gray, thick] (-8.25, 2.5) -- ++ (16.5, 0);
    \draw[dashed, gray, thick] (-8.25,-2.5) -- ++ (16.5, 0);
  \end{scope}


  % ~~~~~~~~~~~~~~~~~~~~~~~~~~~~~~~~~~~~~~~~~~~~~~~~~~~~~~~~~~~~~~~~~~~~~~~~~~~~
  % Series, sequences
  % ~~~~~~~~~~~~~~~~~~~~~~~~~~~~~~~~~~~~~~~~~~~~~~~~~~~~~~~~~~~~~~~~~~~~~~~~~~~~

  % Fibonacci spiral
  \begin{scope}[shift={(\paperwidth*.8,\paperheight*.77)}, scale=1/7.7, xscale=-1]
    \def\inverseSpiral#1{%
      \pgfmathsetmacro\currentVal{int(#1)}
      \ifnum\currentVal>0
        \begin{pgfonlayer}{bg}
          \draw [primaryColor, thick] (0,0) rectangle ++(1,1);
        \end{pgfonlayer}

        % Then draw the square and arc for the current step
        \draw [ultra thick, primaryColor] (0,0) arc (270:360:1);

        % First, recursively call the macro with one less
        \begin{scope}[scale=1.6180339887, rotate=90, shift={(1/1.6180339887,-1/1.6180339887)}]
          \inverseSpiral{#1-1}
        \end{scope}
      \fi
    }

    \inverseSpiral{8}
  \end{scope}

  % sum 1/2^n

  \begin{scope}[shift={(\paperwidth*.065,\paperheight*.625)}, scale=4, draw=primaryColor, very thick]
    \draw (0,0) rectangle            (1,1); % 1
    \draw (0.5,0) --               (0.5,1); % 1/2
    \draw (0.5,0.5) --             (1,0.5); % 1/4
    \draw (0.75,0.5) --           (0.75,1); % 1/8
    \draw (0.75,0.75) --          (1,0.75); % 1/16
    \draw (0.875,0.75) --        (0.875,1); % 1/32
    \draw (0.875,0.875) --       (1,0.875); % 1/64
    \draw (0.9375,0.875) --     (0.9375,1); % 1/128
    \draw (0.9375,0.9375) --    (1,0.9375); % 1/256
    \draw (0.96875,0.9375) --  (0.96875,1); % 1/512
    \draw (0.96875,0.96875) -- (1,0.96875); % 1/1024

    \node at (0.25,0.5) {\scalebox{1.7}{$\sfrac{1}{2}$}};
    \node at (0.75,0.25) {\scalebox{1.2}{$\sfrac{1}{4}$}};
    \node at (0.625,0.75) {$\sfrac{1}{8}$};
    \node at (0.875,0.625) {$\sfrac{1}{16}$};
    \node at (0.8125,0.875) {\scalebox{.8}{$\sfrac{1}{32}$}};
    \node at (0.9375,0.8125) {\scalebox{.8}{$\sfrac{1}{64}$}};
    \node at (0.90625,0.9375) {\scalebox{.3}{$\sfrac{1}{128}$}};
    \node at (0.96875,0.90625) {\scalebox{.2}{$\sfrac{1}{256}$}};
    \node at (0.953125,0.96875) {\scalebox{.1}{$\sfrac{1}{512}$}};
    \node at (0.984375,0.953125) {\scalebox{.05}{$\sfrac{1}{1024}$}};

    % Rest
    \node at (0.984375,0.984375) {\scalebox{.15}{$\dots$}};

  \end{scope}

\end{tikzpicture}

\cleardoublepage


% \begin{titlepage}
%   \includegraphics[width=\paperwidth]{../shared/bme-logo}
% \end{titlepage}

\tableofcontents
\clearpage
\listofdefinitions
\clearpage
\listoftheorems